\documentclass[11pt]{article}
\usepackage[margin=1in]{geometry}
\usepackage[T1]{fontenc}
\usepackage[utf8]{inputenc}
\usepackage{amsmath,amssymb,amsthm}
\usepackage{enumitem}

\title{ICPC World Finals 2019\\K. Traffic Blights}
\author{}
\date{}

\begin{document}
\maketitle

\section*{Problem Summary}

An ideal car appears at a uniformly random time, drives east at speed $1$, and stops forever at the
first red light it reaches. For each light we must compute the probability that it is the first red light hit,
and also the probability that the car never stops.

\section*{Key Observations}

\begin{itemize}[leftmargin=*]
    \item If light $i$ has period $p_i=r_i+g_i$, then the whole system is periodic in the start time with
    period $\operatorname{lcm}(p_1,\dots,p_n)$.
    \item Instead of working modulo that enormous period, fix
    $X = 2^3 \cdot 3^2 \cdot 5 \cdot 7 = 2520$ and split all start times by their residue modulo $X$.
    \item For a fixed residue class $t \bmod X$, light $i$ is then observed only at times
    $t + x_i + kX$. This sequence is periodic with reduced period
    $q_i = p_i / \gcd(p_i, X)$.
    \item By the choice of $X$, every $q_i \le 100$ is a prime power.
    For different primes these reduced periods are coprime; for the same prime they divide one another.
    Therefore the Chinese Remainder Theorem makes the different prime blocks independent.
\end{itemize}

\section*{Algorithm}

\begin{enumerate}[leftmargin=*]
    \item For every light compute its original period $p_i$ and reduced period
    $q_i = p_i / \gcd(p_i, X)$.
    Group all lights with $q_i > 1$ by the prime base of $q_i$.
    For one prime $p$, let $M_p$ be the largest reduced period among those lights.
    \item Fix one residue class $t_0 \in \{0,\dots,X-1\}$.
    For every prime block $p$, maintain which residues modulo $M_p$ are still able to pass all lights
    seen so far. Initially all residues are alive.
    \item Process the lights from west to east.
    \begin{itemize}[leftmargin=*]
        \item If $q_i = 1$, then within this residue class the light has a fixed color.
        It is either always green or always red.
        \item Otherwise, compute for each residue $k \bmod q_i$ whether the car reaches this light at
        a green or red moment.
        Inside the corresponding prime block, count how many currently alive residues modulo $M_p$
        project to a red residue modulo $q_i$.
    \end{itemize}
    \item The product structure implies that, conditioned on passing all previous lights, the residue in
    this prime block is uniformly distributed among the alive residues. Hence the fraction
    $\text{red\_count}/\text{alive\_count}$ is exactly the probability that this light is the first red light
    within the current class $t_0 \bmod X$.
    \item Remove those red residues from the block and continue. The probability mass that survives all
    lights contributes to the final ``passes all lights'' answer.
    \item Average the results over all $X=2520$ residue classes.
\end{enumerate}

\section*{Correctness Proof}

We prove that the algorithm returns the correct probabilities.

\paragraph{Lemma 1.}
Fix a start-time residue class $t_0 \bmod X$.
For light $i$, the color pattern seen at arrival times $t_0 + x_i + kX$ is periodic in $k$ with period
$q_i = p_i / \gcd(p_i, X)$.

\paragraph{Proof.}
Let $d=\gcd(p_i, X)$. Then increasing $k$ by $q_i = p_i/d$ changes the arrival time by
$q_i X = (p_i/d)X$, which is a multiple of $p_i$, so the color repeats.
No smaller period is needed for the algorithm; only periodicity matters here. \qed

\paragraph{Lemma 2.}
Fix $t_0 \bmod X$.
For any set of already processed lights, the set of $k$ values that pass all of them factors as an
independent product over prime blocks of the reduced periods.

\paragraph{Proof.}
Each processed light imposes a condition only on $k \bmod q_i$, and every $q_i$ is a prime power.
Conditions coming from different primes involve coprime moduli, so the Chinese Remainder Theorem
shows they are independent.
Conditions from the same prime all live modulo powers of that prime and therefore can be represented
simultaneously modulo the largest such power in that block. Hence the full feasible set is exactly a
Cartesian product of per-prime feasible residue sets. \qed

\paragraph{Lemma 3.}
While processing a fixed residue class $t_0 \bmod X$, the probability assigned by the algorithm to light
$i$ is exactly the probability that light $i$ is the first red light hit within that residue class.

\paragraph{Proof.}
By Lemma 2, conditioned on having passed all previous lights, each alive residue in the current prime
block is equally likely.
The algorithm counts exactly those alive residues that make light $i$ red, so the ratio
$\text{red\_count}/\text{alive\_count}$ is the conditional probability that light $i$ is red next.
Multiplying by the probability mass that has survived the previous lights gives exactly the probability
that light $i$ is the first red light hit. \qed

\paragraph{Theorem.}
The algorithm outputs the correct probability for every light and for passing all lights.

\paragraph{Proof.}
For each fixed residue class $t_0 \bmod X$, Lemma 3 shows that the algorithm computes the exact
probabilities of first stopping at each light and of surviving all lights.
These residue classes partition all start times into $X$ equally likely classes, so averaging over all
$t_0 \in \{0,\dots,X-1\}$ gives the true overall probabilities. \qed

\section*{Complexity Analysis}

For one residue class $t_0$, each light is processed in $O(q_i + M_p)$ time, where both quantities are
at most $100$.
Thus the total running time is $O(X \cdot n \cdot 100)$, i.e.\ $O(2520 \cdot n \cdot 100)$, which is easily
fast enough for $n \le 500$.
The memory usage is $O(\sum_p M_p)$, at most a few hundred booleans.

\section*{Implementation Notes}

\begin{itemize}[leftmargin=*]
    \item A light is treated as red exactly when the phase is in $[0, r_i)$ and green in $[r_i, p_i)$.
    The boundary has measure zero, so any consistent convention there gives the same probability.
    \item If $q_i = 1$, then within the current class $t_0 \bmod X$ the light's color never changes,
    so it can be handled as a constant-color special case.
\end{itemize}

\end{document}
